\section{Требования физики}
\label{sec:requirements}

Обязательство одно: нижний этаж воспроизводит верхний.
Чтобы было с чем сверять, скажем, что на макроскопическом масштабе установлено прямо, без гипотезы об устройстве вещества.

Возмущение доходит за конечное время.
Если бы влияние было мгновенным, показание прибора в данном месте зависело бы от того,
что в тот же момент происходит сколь угодно далеко, и локальный опыт не сходился бы.
За время~$\caStepTime$ сигнал проходит путь не длиннее~$\speedC\,\caStepTime$:
\begin{equation}
  |\symCapitalDelta\vecX| \le \speedC\,\caStepTime.
  \label{eq:light-cone}
\end{equation}
Любой более глубокий закон обязан сохранить этот конус.

Из \eqref{eq:light-cone} следует локальность.
Далёкое не успевает прийти, поэтому новое состояние в данной точке собирается из поля в её окрестности.
Для электромагнитного поля это уравнения Максвелла:
производная по времени стоит в точке, ротор и дивергенция берут соседние значения,
\begin{align}
  \opNabla\times\fieldE &= -\opPartialT\fieldB, &
  \opNabla\cdot\fieldB &= 0, \label{eq:maxwell-hom}\\
  \opNabla\times\fieldB
    &= \constMuZero\fieldJ + \constMuZero\constEpsZero\,\opPartialT\fieldE, &
  \opNabla\cdot\fieldE &= \chargeRho/\constEpsZero. \label{eq:maxwell-inh}
\end{align}
Характеристики этой системы лежат на конусе \eqref{eq:light-cone}.
Та же схема у гидродинамики.
Для идеальной жидкости скорость в точке меняется от давления и скорости рядом:
\begin{equation}
  \opPartialT\fieldV + (\fieldV\cdot\opNabla)\fieldV
  = -\frac{1}{\massDensity}\opNabla \pressureP.
  \label{eq:euler}
\end{equation}
В механике сплошной среды справа стоит дивергенция напряжений,
$\opPartialK\stressSigmaIk$, снова из окрестности той же точки.

Механическое состояние исчерпывается координатами и скоростями.
Положение точки есть $\fieldR(t)$, скорость --- его производная,
\begin{equation}
  \fieldV = \dot{\fieldR}.
  \label{eq:velocity}
\end{equation}
Уравнения Ньютона связывают ускорение с этим состоянием:
\begin{equation}
  \massBody\ddot{\fieldR} = \fieldF(\fieldR,\dot{\fieldR}).
  \label{eq:newton}
\end{equation}
Это уравнение второго порядка.
Задание $\fieldR$ и~$\dot{\fieldR}$ в один момент восстанавливает их и в предыдущий, и в следующий.
Старшей производной среди начальных данных в \eqref{eq:newton} нет:
закон, которому она понадобилась бы, описывал бы другой мир.

Тепловое равновесие добавляет два запрета, тоже измеренных.
В адиабатически замкнутой системе энтропия не убывает,
\begin{equation}
  \mathrm{d}\entropyS \ge 0.
  \label{eq:second-law}
\end{equation}
При $\tempT\to 0$ энтропия равновесной системы перестаёт зависеть от параметров равновесия~$\symEll$
и в шкале Планка обращается в нуль:
\begin{equation}
  \Bigl(\frac{\partial \entropyS}{\partial \symEll}\Bigr)_{\tempT}
    \xrightarrow[\tempT\to 0]{} 0,
  \qquad
  \entropyS \xrightarrow[\tempT\to 0]{} 0.
  \label{eq:third-law}
\end{equation}
Система садится в одно основное состояние.
Закон, у которого при нулевой температуре оставалось бы макроскопически много
равноправных состояний, \eqref{eq:third-law} противоречит.

Вакуум, насколько его видит свет, не выделяет направления:
скорость одна и та же вдоль любой оси.
Вместе с \eqref{eq:light-cone} это интервал
\begin{equation}
  \mathrm{d}s^2 = \speedC^2\,\mathrm{d}t^2 - \mathrm{d}\vecX^2,
  \label{eq:interval}
\end{equation}
один и тот же во всех инерциальных системах.
Заряд в замкнутой системе сохраняется.
Это уравнение непрерывности
\begin{equation}
  \opPartialT\chargeRho + \opNabla\cdot\fieldJ = 0,
  \label{eq:continuity}
\end{equation}
а на узле проводников --- закон Кирхгофа, $\sum_k I_k = 0$.
Не свойство, приписанное заранее микроскопическому полю.

На этом макроскопическая сверка кончается.
Спин, комплексная амплитуда, целочисленность устойчивого заряда
и вопрос о том, откуда берутся безразмерные числа, здесь ещё не видны.
Они появятся на том этаже, где без них уже проверенный опыт перестаёт получаться,
и каждое окажется тем же конусом, тем же сохранением или той же обратимостью, дочитанными до этого этажа.

